\documentclass[10pt,letterpaper]{article}
\usepackage{technical-user-manual}
\setDocTitle{Chapter 35: Follow Through}
\setDocSubtitle{Book-aligned guided reading and application}
\setDocVersion{1.0}
\setDocStatus{Educational / Draft}
\setDocOwner{Game Development}
\setDocDate{2026-10-30}
\setDocSummary{Targets Chapter 35 of Timing for Animation, third edition.}
\begin{document}
\makedoctitle
\makedocsummary

\section{Reading map and provenance}
\textbf{Source reading:} Whitaker, Halas, and Sito (2020), Chapter 35, ``Follow Through.``\\
\textbf{Domain progression:} Spacing, holds, anticipation, overlap, weight, and cycles.
\begin{itemize}
\item Source-aligned: read the named chapter; identify its terminology, illustrations, production assumptions, and the reason its timing choices are made.
\item Supplementary engineering: the implementation assignment below applies those artistic ideas to a modern programmable animation system; it is not claimed as a recipe from the book.
\end{itemize}
\section{Target competency}
Distinguish primary movement from follow-through. Supply a short explanation and demonstrate the behavior in a repeatable shot.
\section{Guided analysis}
\begin{itemize}
\item Which specific visual or production problem does the chapter address in follow through?
\item What changes if the frame interval or available animation budget changes while the intended action remains the same?
\item Which observation or pose sequence would distinguish a successful result from an unsuccessful one?
\item Which assumptions in the chapter are specific to a production medium, and which transfer to real-time games?
\end{itemize}
\section{Engineering translation (supplementary)}
\textbf{Artifact:} settle offset chart. Record fps as a rational value, duration in seconds, frame boundaries, and whether transforms are evaluated from elapsed time or counted exposures.
\textbf{Exercise:} Quantify delayed accessory movement after a sudden stop. Prepare a baseline, one controlled timing variation, and a comparison. Keep pose endpoints and camera setup fixed unless the exercise is specifically about those variables.
\section{Evidence and self-assessment}
\begin{itemize}
\item \textbf{Reproduction:} include shot identifier, original reference observations, frames/seconds and any relevant timecode convention.
\item \textbf{Control:} document the changed variable, why the change was selected, and the expected perceptual effect.
\item \textbf{Assessment:} capture representative key poses, a playback review at the intended display rate, and at least one failure mode.
\item \textbf{Reflection:} distinguish artistic intent, physical plausibility, and real-time implementation constraints.
\end{itemize}
\section{Cross-reference}
Continue to the timing technical reference, the animation review scorecard, and the curriculum prerequisite map in this suite. This guide is independently authored and is not a substitute for the chapter's figures or explanations.
\end{document}
